\chapter{Blood Culture Collection}

\section*{Overview}
Blood culture collection is the sterile withdrawal of blood for culture to detect bacteremia, fungemia, and bloodstream infection. Specimens may be obtained from a peripheral venepuncture or from an existing central venous catheter.
\section*{Alternative Names}
Blood cultures; bacteremia testing; blood sampling for culture; central line blood culture
\section*{Principle}
A small volume of blood is inoculated into culture bottles containing media that support bacterial and fungal growth. Because blood is normally sterile, any growth indicates a true infection once skin contaminants are excluded, and the organism's antibiotic susceptibility can guide therapy.
\section*{History}
The blood culture was introduced in the late 19th century and became a cornerstone of the diagnosis of sepsis. Modern automated culture systems continuously monitor bottles for growth and have greatly improved speed and yield.
\section*{Why It Is Done}
Ordered for suspected sepsis, fever of unknown origin, endocarditis, pneumonia with bacteremia, and monitoring of patients with central venous catheters at risk of catheter-related bloodstream infection.
\section*{Is This a Primary or Secondary Test or Procedure}
Primary diagnostic procedure for bloodstream infection and a key driver of antimicrobial therapy.
\section*{How It Is Performed}
The skin is disinfected with chlorhexidine and allowed to dry. Blood is drawn from a peripheral vein or through a central catheter after disinfecting the access hub, and each bottle is filled with the recommended volume. Two to three sets are typically collected from different sites, and the bottles are labelled with the site and time.
\section*{Preparation}
No patient preparation is required, but cultures should be drawn before antibiotics when possible, using strict aseptic technique. Ideal volumes are collected from separate sites rather than a single draw.
\section*{Contraindications and Precautions}
There are no absolute contraindications; precautions include drawing from a central catheter only when peripheral access is not possible, and avoiding cultures drawn through an obviously infected or freshly accessed line.
\section*{Risks}
Pain and bruising at the venepuncture site, bleeding, and contamination of the sample with skin flora that produces false-positive results.
\section*{Aftercare and Recovery}
Bottles are transported promptly to the laboratory and incubated in automated systems, typically for up to five days. Positive growth triggers Gram stain and identification, with susceptibility testing reported to guide therapy.
\section*{Outcome and Follow-up}
A positive culture identifies the causative organism, and susceptibility results guide targeted antibiotics. Growth in multiple sets from different sites strongly supports true bacteremia, whereas growth in a single set may represent contamination.
\section*{Expected Results}
No growth of organisms from the blood cultures after the full incubation period.
\section*{Follow-up}
Patients with positive cultures may need repeat cultures to confirm clearance, echocardiography for suspected endocarditis, and imaging or source control when a focus of infection is identified.
\section*{When to Seek Medical Care}
Follow the procedure team's specific instructions. Seek urgent medical assessment for severe or worsening pain, uncontrolled bleeding, fever or signs of infection, trouble breathing, chest pain, fainting, new weakness or confusion, inability to urinate, or any rapidly worsening symptom. The appropriate warning signs depend on the procedure and the patient's medical condition.

\section*{References}
\begin{enumerate}
  \item MedlinePlus. Blood culture. U.S. National Library of Medicine; 2025.
  \item Current Medical Diagnosis and Treatment. McGraw-Hill; 2025.
\end{enumerate}

\clearpage
